\begin{abstract}
Retrieval-Augmented Generation (RAG) has become the dominant paradigm for grounding large language model outputs in external knowledge. Yet practitioners face a vast configuration space spanning chunking strategies, embedding models, retrieval methods, rerankers, prompt designs, and generation models---with limited guidance on how these components interact. We present \toolname{}, a modular benchmarking framework that attaches to existing RAG pipelines as a non-invasive evaluation layer. It systematically explores the full configuration space via cartesian product sweeps, replacing manual tuning with reproducible experiments. Every configuration is tracked end-to-end through MLflow and evaluated through a hybrid metric scheme combining LLM-based quality assessment (RAGAS), established information retrieval measures (nDCG, Hit@k, Recall@k), natural language generation metrics (BERTScore, METEOR, ROUGE-L, BLEU), and runtime performance indicators (latency, throughput). Pre-configured dataset adapters cover diverse QA domains and are extensible to custom data. We validate the framework on 182 configurations across seven generation models spanning local, cloud-hosted, and API deployments on SQuAD, demonstrating non-obvious interactions: MMR retrieval improves nDCG@5 by up to 21\% over similarity search, detailed prompt templates consistently raise faithfulness by 6--11 percentage points, and larger generation models yield higher answer quality at the cost of increased latency. The framework enables systematic, comparable evaluation of RAG components with minimal integration effort.

\keywords{Retrieval-Augmented Generation \and Benchmarking Framework \and RAG Evaluation \and Configuration Analysis}
\end{abstract}
